\documentclass[11pt]{article}
\usepackage[margin=1in]{geometry}
\usepackage{amsmath, amssymb, amsthm}
\usepackage{mathtools}
\usepackage{hyperref}

\newtheorem{theorem}{Theorem}
\newtheorem{proposition}{Proposition}
\newtheorem{lemma}{Lemma}
\newtheorem{corollary}{Corollary}

\begin{document}

\begin{theorem}[Equivalence of exact privacy curves for additive symmetric noise]
Let
\[
M(x)=A(x)+Z,
\]
where \(A:\mathcal X\to \mathbb R^d\) is deterministic and \(Z\sim \nu\), with \(\nu\) a centered symmetric probability distribution on \(\mathbb R^d\), i.e.
\[
\nu(B)=\nu(-B)\qquad \text{for every Borel set }B\subseteq \mathbb R^d.
\]
Fix neighboring datasets \(x,x'\), and define
\[
a:=A(x),\qquad a':=A(x'),
\]
\[
P:=\mathcal L(a+Z),\qquad Q:=\mathcal L(a'+Z).
\]

Assume \(P\ll Q\) (equivalently \(Q\ll P\) in this setting). Let
\[
L:=\frac{dP}{dQ}.
\]

Then the following data determine one another:
\begin{enumerate}
\item the full exact R\'enyi-DP curve \(\{D_\alpha(P\|Q):\alpha>1\}\);
\item the full exact \(f\)-DP curve;
\item the exact \((\varepsilon,\delta)\)-DP curve for \(\varepsilon\ge 0\).
\end{enumerate}

Equivalently, in this symmetric additive-noise setting, the one-sided exact privacy profile \(\{\delta(\varepsilon):\varepsilon\ge 0\}\) already determines the full two-sided profile \(\{\delta(\varepsilon):\varepsilon\in\mathbb R\}\), and hence determines the R\'enyi-DP and \(f\)-DP curves. The converse implications also hold.
\end{theorem}

\begin{proof}
We proceed in several steps.

\medskip
\noindent
\textbf{Step 1: Reflection symmetry exchanges \(P\) and \(Q\).}

Define the midpoint
\[
m:=\frac{a+a'}{2},
\]
and the reflection
\[
T(y):=2m-y=a+a'-y,\qquad y\in\mathbb R^d.
\]
Then \(T\) is a measurable involution: \(T\circ T=\mathrm{id}\).

We claim that
\[
T_\# P = Q,\qquad T_\# Q = P.
\]
Indeed, if \(Y=a+Z\sim P\), then
\[
T(Y)=a+a'-(a+Z)=a'-Z.
\]
Since \(Z\stackrel d= -Z\) by symmetry of \(\nu\), we have
\[
a'-Z \stackrel d= a'+Z \sim Q.
\]
Hence \(T_\#P=Q\). By symmetry the reverse identity \(T_\#Q=P\) also holds.

\medskip
\noindent
\textbf{Step 2: The privacy profile is symmetric under swapping \(P\) and \(Q\).}

Recall the exact privacy profile
\[
\delta_{P,Q}(\varepsilon):=\sup_{E}\bigl(P(E)-e^\varepsilon Q(E)\bigr),\qquad \varepsilon\in\mathbb R.
\]
For any measurable \(E\),
\[
P(E)-e^\varepsilon Q(E)
=Q(T^{-1}E)-e^\varepsilon P(T^{-1}E).
\]
Taking the supremum over \(E\), and using that \(T^{-1}\) is a bijection on measurable sets, yields
\[
\delta_{P,Q}(\varepsilon)=\delta_{Q,P}(\varepsilon)\qquad \forall \varepsilon\in\mathbb R.
\]
Thus the privacy profile is invariant under exchanging \(P\) and \(Q\).

\medskip
\noindent
\textbf{Step 3: Negative \(\varepsilon\) is determined by positive \(\varepsilon\).}

For arbitrary probability measures \(P,Q\), one has the identity
\[
\delta_{P,Q}(-\varepsilon)
=
1-e^{-\varepsilon}+e^{-\varepsilon}\delta_{Q,P}(\varepsilon),
\qquad \varepsilon\in\mathbb R.
\]
Indeed,
\[
\delta_{P,Q}(-\varepsilon)
=
\sup_E \bigl(P(E)-e^{-\varepsilon}Q(E)\bigr).
\]
Writing \(F=E^c\), we get
\begin{align*}
P(E)-e^{-\varepsilon}Q(E)
&=1-P(F)-e^{-\varepsilon}(1-Q(F))\\
&=1-e^{-\varepsilon}+\bigl(e^{-\varepsilon}Q(F)-P(F)\bigr)\\
&=1-e^{-\varepsilon}+e^{-\varepsilon}\bigl(Q(F)-e^\varepsilon P(F)\bigr).
\end{align*}
Taking the supremum over \(E\), equivalently over \(F\), proves the identity.

Now apply Step 2:
\[
\delta_{Q,P}(\varepsilon)=\delta_{P,Q}(\varepsilon).
\]
Therefore
\[
\boxed{
\delta(-\varepsilon)=1-e^{-\varepsilon}+e^{-\varepsilon}\delta(\varepsilon),
\qquad \varepsilon\ge 0.
}
\]
Hence knowledge of \(\delta(\varepsilon)\) for all \(\varepsilon\ge 0\) uniquely determines \(\delta(\varepsilon)\) for all \(\varepsilon<0\). So the one-sided privacy profile determines the full two-sided profile.

\medskip
\noindent
\textbf{Step 4: The full two-sided privacy profile determines the law of the likelihood ratio.}

Let
\[
L:=\frac{dP}{dQ}.
\]
As usual,
\[
\delta(\varepsilon)=\mathbb E_Q[(L-e^\varepsilon)_+],\qquad \varepsilon\in\mathbb R.
\]
Equivalently, if
\[
g(u):=\mathbb E_Q[(L-u)_+],\qquad u>0,
\]
then
\[
\delta(\varepsilon)=g(e^\varepsilon).
\]
For every \(u>0\), \(g\) is convex and has one-sided derivatives
\[
g'_-(u)=-Q(L\ge u),\qquad g'_+(u)=-Q(L>u).
\]
Thus knowing \(g(u)\) for all \(u>0\) determines the tail probabilities of \(L\), hence the law of \(L\) under \(Q\).

Since Step 3 shows that the one-sided profile \(\{\delta(\varepsilon):\varepsilon\ge 0\}\) determines the full profile \(\{\delta(\varepsilon):\varepsilon\in\mathbb R\}\), it follows that in the present symmetric setting, the one-sided profile determines the law of \(L\).

\medskip
\noindent
\textbf{Step 5: The law of \(L\) determines the full exact R\'enyi-DP curve and the full exact \(f\)-DP curve.}

For every \(\alpha>1\),
\[
D_\alpha(P\|Q)=\frac{1}{\alpha-1}\log \mathbb E_Q[L^\alpha].
\]
Hence the law of \(L\) determines the full R\'enyi-DP curve.

Likewise, by the Neyman--Pearson lemma, the exact \(f\)-DP / tradeoff curve is determined by threshold tests in \(L\), hence is determined by the law of \(L\).

Therefore the one-sided exact privacy profile determines both the full exact R\'enyi-DP curve and the full exact \(f\)-DP curve.

\medskip
\noindent
\textbf{Step 6: Converse implications.}

The full exact R\'enyi-DP curve determines the law of \(L\): indeed, for \(\alpha>1\),
\[
\mathbb E_Q[L^\alpha]=e^{(\alpha-1)D_\alpha(P\|Q)},
\]
so the moment generating function of \(\log L\) is known on \((1,\infty)\), which determines its distribution uniquely.

The full exact \(f\)-DP curve also determines the law of \(L\), by Neyman--Pearson.

Once the law of \(L\) is known, the exact privacy profile is given by
\[
\delta(\varepsilon)=\mathbb E_Q[(L-e^\varepsilon)_+].
\]
In particular, this determines the one-sided profile \(\varepsilon\ge 0\).

Thus all three objects determine one another.
\end{proof}


\end{document}